\scanpage{5}
\mspage{2}
de plus, on vérifie fibre par fibre que
\[
R^1p_*\handwritten{\mathscr O}(ne)=0\quad\text{pour}\quad n>0\qquad .
\]

Il en résulte, par Riemann-Roch, que pour $n>0$, le faisceau $p_*\handwritten{\mathscr O}(ne)$ est localement libre de rang $n$, et que sa formation commute à tout changement de base. La suite exacte de cohomologie fournit de plus une suite exacte
\[
0\handwritten{\longrightarrow}p_*\handwritten{\mathscr O}(ne)\handwritten{\longrightarrow}
p_*\handwritten{\mathscr O}((n+1)e)\handwritten{\longrightarrow}\omega^{\otimes-(n+1)}
\handwritten{\longrightarrow}0\quad(n>0)\quad .
\]
Par ailleurs, $\handwritten{\mathscr O}(e)$ étant localement libre de rang 1 , on vérifie fibre par fibre, compte tenu \unsure{dé|de} (SGA 1 X 1.2), que
\[
\handwritten{\mathscr O}_S\handwritten{\xrightarrow{\sim}}p_*\handwritten{\mathscr O}_E
\handwritten{\xrightarrow{\sim}}p_*\handwritten{\mathscr O}(e)\qquad .
\]
En d'autres termes, les $p_*\handwritten{\mathscr O}(me)$ pour $1\leq m\leq n$ définissent une filtration de $p_*\handwritten{\mathscr O}(ne)$, de gradué associé
\[
\mathrm{Gr}\ p_*\handwritten{\mathscr O}(ne)=\sum_{\substack{i=0\\i\ne1}}^n\omega^{\otimes-i}\qquad .\tag*{(1.1)}
\]

Le morphisme $p$ est d'intersection complète relative ; il en résulte que $Rp^!\handwritten{\mathscr O}_S\in D^+(E)$ \handwritten{\tikz[baseline=-.2ex,x=1ex,y=1ex]{\draw (0,0)--(.28,1.22) .. controls (.13,1.48) and (.51,1.55) .. (.47,1.35) .. controls (.44,1.22) and (.30,1.22) .. (.28,1.22);}}\margin{\handwritten{\uline{réf}}} a un seul groupe de cohomologie non nulle, $R^1p^!\handwritten{\mathscr O}_S$, et que ce dernier est un faisceau inversible sur $E$ ; on l'appelle le faisceau des différentielles régulières sur $E$ et on le désigne par $\Omega^{\mathrm{reg}}_{E/S}$. Là où $p$ est lisse, on a $\Omega^{\mathrm{reg}}_{E/S}=\Omega^1_{E/S}$. La dualité montre que $R^1p_*\Omega^{\mathrm{reg}}_{E/S}$ est isomorphe à $\handwritten{\mathscr O}$, $p_*\Omega^{\mathrm{reg}}_{E/S}$ est donc\handwritten{\tikz[baseline=-.2ex,x=1ex,y=1ex]{\draw (0,0)--(.28,1.22) .. controls (.13,1.48) and (.51,1.55) .. (.47,1.35) .. controls (.44,1.22) and (.30,1.22) .. (.28,1.22);}}\margin{\handwritten{\uline{réf}}} localement libre, de formation compatible à tout changement de base et, raisonnant fibre par fibre, on en déduit que

% archive sheet number: 4
